\documentclass[a4paper,12pt]{article}
\usepackage{docsTemplate}
\usepackage{eurosym}
\usepackage{array}
\usepackage{float}
\newcolumntype{M}[1]{>{\centering\arraybackslash}m{#1}}

\setTitle{Piano di Qualifica}
\setAuthors{Andrea Masiero \\ & Bianca Zaghetto \\ & Mattia Oliva Medin \\ & Sara Gioia Fichera}
\setVerificators{ Mattia Oliva Medin \\ & Filippo Compagno \\ & Nenad Radulovic \\ & Francesco Balestro }
\setApprovation{Nenad Radulovic}
\setVersion{1.2}
\setType{Esterno}
\setDestination{Tullio Vardanega \\ & Riccardo Cardin \\ & Team Three Technologies}

\begin{document}
    \include{docTitlepage}

    \section*{Registro delle modifiche} {
        \begin{table}[h!]
            \rowcolors{2}{lightgray}{white}
                \begin{tabularx}{\textwidth}{|l|l|l|l|L|L|}
                \hline
                \textbf{Vers.} & \textbf{Data} & \textbf{Autore} & \textbf{Ruolo} & \textbf{Verifica} & \textbf{Oggetto} \\
                \hline
                \textbf{1.2} & 14/04/2026 & Bianca Zaghetto & Programmatore & Filippo Compagno & Inserimento grafici\\
                \hline
                \textbf{1.1} & 13/04/2026 & Bianca Zaghetto & Programmatore & Francesco Balestro & Aggiornamento per sprint 13\\
                \hline
                \textbf{1.0} & 18/03/2026 & Nenad Radulovic & Responsabile & - & Approvazione e rilascio\\
                \hline
                \textbf{0.10} & 17/03/2026 & Mattia Oliva Medin & Progettista & Andrea Masiero & Aggiornamento grafici per sprint 9\\
                \hline
                \textbf{0.9} & 26/02/2026 & Francesco Balestro & Progettista & Sara Gioia Fichera & Inserimento grafici\\
                \hline
                \textbf{0.8} & 3/02/2026 & Sara Gioia Fichera & Programmatore & Francesco Balestro & Modifica test di Sistema\\
                \hline
                \textbf{0.7} & 12/01/2026 & Sara Gioia Fichera & Amministratore & Mattia Oliva Medin & Correzione errori e completamento metriche \\
                \hline
                \textbf{0.6} & 7/01/2026 & Andrea Masiero & Analista & Filippo Compagno & Correzione test \\
                \hline
                \textbf{0.5} & 3/01/2026 & Sara Gioia Fichera & Responsabile & Mattia Oliva Medin & Aggiunta sezione test di Sistema \\
                \hline
                \textbf{0.4} & 3/01/2026 & Bianca Zaghetto & Amministratore & Filippo Compagno & Stesura checklist \\
                \hline
                \textbf{0.3} & 2/01/2026 & Bianca Zaghetto & Amministratore & Filippo Compagno & Aggiunta sezione test di Accettazione \\
                \hline
                \textbf{0.2} & 19/12/2025 & Mattia Oliva Medin & Amministratore & Nenad Radulovic & Stesura metriche\\
                \hline
                \textbf{0.1} & 18/12/2025 & Sara Gioia Fichera & Analista & Nenad Radulovic & Struttura del documento \\
                \hline
            \end{tabularx}
        \end{table}
    }
    \newpage
    
    \tableofcontents
    \newpage

    \listoffigures
    \newpage

    \listoftables
    \newpage
    
    \section{Introduzione} {
        \subsection{Scopo del documento}
            Il presente documento ha lo scopo di definire le strategie, le metodologie e le metriche adottate per garantire la qualità del prodotto software e l'efficienza dei processi di sviluppo.

        \subsection{Glossario}
            Si rimanda al documento \textit{Glossario} per chiarire i termini nel dominio del progetto che possono risultare ambigui o di difficile comprensione. I vocabili presenti nel glossario sono stati segnalati seguendo la notazione:
            \begin{center}
                parola\textsubscript{G}
            \end{center}
        \subsection{Riferimenti}
            \subsubsection{Riferimenti normativi}
                \begin{itemize}
                    \item \textit{Norme di progetto v1.0}
                    \begin{itemize}
                        \item \url{https://team-three-technologies.github.io/Documenti/2%20-%20PB/Documenti%20Interni/Norme_di_Progetto_v1_0.pdf}
                        \item[] (Ultimo accesso: 2026/04/11)
                    \end{itemize}
                    \item Capitolato C3
                    \begin{itemize}
                        \item \url{https://www.math.unipd.it/~tullio/IS-1/2025/Progetto/C3.pdf}
                        \item[] (Ultimo accesso: 2026/01/11)
                    \end{itemize}
                    \item Regolamento del Progetto Didattico
                    \begin{itemize}
                        \item \url{https://www.math.unipd.it/~tullio/IS-1/2025/Dispense/PD1.pdf}
                        \item[] (Ultimo accesso: 2026/04/05)
                    \end{itemize}
                \end{itemize}
            \subsubsection{Riferimenti informativi}
                \begin{itemize}
                    \item Verbali interni
                    \item Verbali esterni
                    \item \textit{Analisi dei Requisiti v1.1}
                    \begin{itemize}
                        \item \url{https://team-three-technologies.github.io/Documenti/2%20-%20PB/Documenti%20esterni/Analisi_dei_Requisiti_v1_1.pdf}
                        \item[] (Ultimo accesso: 2026/04/11)
                    \end{itemize}
                    \item \textit{Piano di Progetto v1.4}
                    \begin{itemize}
                        \item \url{https://team-three-technologies.github.io/Documenti/2%20-%20PB/Documenti%20esterni/Piano_di_Progetto_v1_4.pdf}
                        \item[] (Ultimo accesso: 2026/04/11)
                    \end{itemize}
                    \item Slide del corso: T7 - Qualità di prodotto
                    \begin{itemize}
                        \item \url{https://www.math.unipd.it/~tullio/IS-1/2025/Dispense/T07.pdf}
                         \item[] (Ultimo accesso: 2026/04/06)
                    \end{itemize}
                    \item Slide del corso: T8 - Qualità di processo
                    \begin{itemize}
                        \item \url{https://www.math.unipd.it/~tullio/IS-1/2025/Dispense/T08.pdf}
                        \item[] (Ultimo accesso: 2026/04/11)
                    \end{itemize}
                    \item Standard ISO/IEC 12207
                    \begin{itemize}
                        \item \url{https://www.math.unipd.it/~tullio/IS-1/2009/Approfondimenti/ISO_12207-1995.pdf}
                        \item[] (Ultimo accesso: 2026/03/03)
                    \end{itemize}
                    \item \textit{Glossario v0.5}
                    \begin{itemize}
                        \item \url{https://team-three-technologies.github.io/Documenti/1%20-%20PB/Documenti%20interni/Glossario_v0_5.pdf}
                        \item[] (Ultimo accesso: 2026/04/06)
                    \end{itemize}
                \end{itemize}
    }

    \newpage

    \section{Obiettivi di qualità}
    La presente sezione definisce gli obiettivi di qualità del progetto, specificando per ciascuna metrica definita all'interno del documento di \textit{Norme di Progetto} (\S5) i range di accettazione e i valori ideali prefissati dal gruppo. La classificazione segue la distinzione tra:
    \begin{itemize}
        \item Metriche di processo
        \item Metriche di prodotto
    \end{itemize}
    
        \subsection{Qualità di processo}
            In questa sezione le metriche vengono riorganizzate seguendo questa divisione: 
            \begin{itemize}
                \item Processi Primari: Obiettivi legati alla fornitura, allo sviluppo e al controllo dei costi/tempi, nonché alla codifica.
                \item Processi di Supporto: Obiettivi trasversali legati alla documentazione, alla verifica e alla gestione della configurazione.
                \item Processi Organizzativi: Obiettivi legati alla gestione delle risorse e al miglioramento continuo.
            \end{itemize}
            
            \subsubsection{Processi Primari}
                \paragraph{Fornitura} \mbox{}
                    \begin{longtable}{|M{0.15\linewidth}|M{0.35\linewidth}|M{0.2\linewidth}|M{0.2\linewidth}|}

                        \caption{Valori metriche processo di fornitura} \\
                        \hline
                        \rowcolor{gray!40}
                        \textbf{ID} & \textbf{Nome Metrica} & \textbf{Valore Accettabile} & \textbf{Valore Ideale} \\
                        \hline
                        \endhead

                        MPC01 & Metriche soddisfatte & $\ge 75\%$ & $100\%$ \\
                        \hline

                        MPC02 & Earned Value & $\ge 0$ & $\le EAC$ \\
                        \hline
                        
                        MPC03 & Planned Value & $\ge 0$ & $\le BAC$ \\
                        \hline
                        
                        MPC04 & Actual Cost & $\ge 0$ & $\le EAC$\\
                        \hline
                        
                        MPC05 & Estimated to Completion & $\ge 0$ & $\le EAC$ \\
                        \hline

                        MPC06 & Estimated at Completion & $\le BAC + 5\%$ & $\le BAC$ \\
                        \hline

                        MPC07 & Cost Variance & $\ge -5\%$ & $\ge 0\%$\\
                        \hline

                        MPC08 & Schedule Variance & $\ge -5\%$ & $\ge 0\%$\\
                        \hline

                        MPC09 & Cost Performance Variance & $\ge 0.90$ & 1 \\
                        \hline

                    \end{longtable}

                    \paragraph{Sviluppo} \mbox{}

                        \begin{longtable}{|M{0.15\linewidth}|M{0.35\linewidth}|M{0.2\linewidth}|M{0.2\linewidth}|}

                        \caption{Valori metriche processo di sviluppo} \\
                        \hline
                        \rowcolor{gray!40}
                        \textbf{ID} & \textbf{Nome Metrica} & \textbf{Valore Accettabile} & \textbf{Valore Ideale} \\
                        \hline
                        \endhead

                        MPC10 & Requirements Stability Index & $\ge 60\%$ & $100\%$ \\
                        \hline

                        \end{longtable}

            \subsubsection{Processi di Supporto}
                \paragraph{Documentazione} \mbox{}
                    \begin{longtable}{|M{0.15\linewidth}|M{0.35\linewidth}|M{0.2\linewidth}|M{0.2\linewidth}|}

                        \caption{Valori metriche processo di documentazione} \\
                        \hline
                        \rowcolor{gray!40}
                        \textbf{ID} & \textbf{Nome Metrica} & \textbf{Valore Accettabile} & \textbf{Valore Ideale} \\
                        \hline
                        \endhead

                        MPC11 & Indice Gulpease & $\ge 40\%$ & $\ge 60\%$ \\
                        \hline

                        MPC12 & Correttezza Ortografica & $0$ & $0$ \\
                        \hline

                    \end{longtable}

                \paragraph{Gestione della qualità} \mbox{}
                    \begin{longtable}{|M{0.15\linewidth}|M{0.35\linewidth}|M{0.2\linewidth}|M{0.2\linewidth}|}

                        \caption{Valori metriche processo di gestione della qualità} \\
                        \hline
                        \rowcolor{gray!40}
                        \textbf{ID} & \textbf{Nome Metrica} & \textbf{Valore Accettabile} & \textbf{Valore Ideale} \\
                        \hline
                        \endhead

                        MPC13 & Metriche non soddisfatte & $\le 2$ & $0$ \\
                        \hline

                        MPC14 & Test Pass Rate & $-$ & $-$ \\
                        \hline

                    \end{longtable}


            \subsubsection{Processi organizzativi}
                \paragraph{Gestione} \mbox{}
                    \begin{longtable}{|M{0.15\linewidth}|M{0.35\linewidth}|M{0.2\linewidth}|M{0.2\linewidth}|}

                        \caption{Valori metriche processi di gestione} \\
                        \hline
                        \rowcolor{gray!40}
                        \textbf{ID} & \textbf{Nome Metrica} & \textbf{Valore Accettabile} & \textbf{Valore Ideale} \\
                        \hline
                        \endhead

                        MPC15 & Rischi non calcolati & $\le 2$ & $0$ \\
                        \hline

                        MPC16 & Efficienza Temporale & $200$ & $100$ \\
                        \hline

                    \end{longtable}



        \subsection{Qualità di prodotto}
            \subsubsection{Funzionalità}
                \begin{longtable}{|M{0.15\linewidth}|M{0.35\linewidth}|M{0.2\linewidth}|M{0.2\linewidth}|}

                    \caption{Valori metriche funzionalità di prodotto} \\
                    \hline
                    \rowcolor{gray!40}
                    \textbf{ID} & \textbf{Nome Metrica} & \textbf{Valore Accettabile} & \textbf{Valore Ideale} \\
                    \hline
                    \endhead

                    MPD01 & Copertura dei requisiti obbligatori & $100\%$ & $100\%$ \\
                    \hline

                    MPD02 & Copertura dei requisiti desiderabili & $0\%$ & $100\%$ \\
                    \hline

                    MPD03 & Copertura dei requisiti opzionali & $0\%$ & $100\%$ \\
                    \hline

                \end{longtable}

            \subsubsection{Affidabilità}
            \begin{longtable}{|M{0.15\linewidth}|M{0.35\linewidth}|M{0.2\linewidth}|M{0.2\linewidth}|}

                \caption{Valori metriche affidabilità di prodotto} \\
                \hline
                \rowcolor{gray!40}
                \textbf{ID} & \textbf{Nome Metrica} & \textbf{Valore Accettabile} & \textbf{Valore Ideale} \\
                \hline
                \endhead

                    MPD04 & Code coverage & $\ge80\%$ & $\ge90\%$ \\
                    \hline

                    MPD05 & Branch Coverage & $\ge60\%$ & $\ge80\%$ \\
                    \hline

                    MPD06 & Statement Coverage & $\ge60\%$ & $\ge80\%$ \\
                    \hline

            \end{longtable}



            \subsubsection{Usabilità}
            \begin{longtable}{|M{0.15\linewidth}|M{0.35\linewidth}|M{0.2\linewidth}|M{0.2\linewidth}|}

                \caption{Valori metriche usabilità di prodotto} \\
                \hline
                \rowcolor{gray!40}
                \textbf{ID} & \textbf{Nome Metrica} & \textbf{Valore Accettabile} & \textbf{Valore Ideale} \\
                \hline
                \endhead

                    MPD07 & Error rate & $\le2$ & $0$ \\
                    \hline

                    MPD08 & Tempo di apprendimento & $ - $ & $ - $ \\
                    \hline

            \end{longtable}

            \subsubsection{Efficienza}
            \begin{longtable}{|M{0.15\linewidth}|M{0.35\linewidth}|M{0.2\linewidth}|M{0.2\linewidth}|}

                \caption{Valori metriche efficienza di prodotto} \\
                \hline
                \rowcolor{gray!40}
                \textbf{ID} & \textbf{Nome Metrica} & \textbf{Valore Accettabile} & \textbf{Valore Ideale} \\
                \hline
                \endhead

                    MPD09 & Response time & $\le 1sec $ & $\le 1sec $  \\
                    \hline

            \end{longtable}


            \subsubsection{Manutenibilità}
            \begin{longtable}{|M{1.5cm}|M{5cm}|M{2.5cm}|M{2.5cm}|}

                \caption{Valori metriche manutenibilità di prodotto} \\
                \hline
                \rowcolor{gray!40}
                \textbf{ID} & \textbf{Nome Metrica} & \textbf{Valore Accettabile} & \textbf{Valore Ideale} \\
                \hline
                \endhead

                    MPD10 & Accoppiamento tra classi & $\le4$ & $\le2$ \\
                    \hline

                    MPD11 & Complessità ciclomatica & $\le15$ & $\le10$ \\
                    \hline

            \end{longtable}
            
\newpage
    \section{Testing}

    Questa sezione descrive la strategia di verifica e validazione adottata dal gruppo per garantire la qualità del prodotto software.
    L'attività di testing è concepita come un processo continuo e fondamentale per rilevare tempestivamente eventuali difetti, assicurando che il Sistema soddisfi tutti i requisiti concordati.

        \subsection{Test di Unità}

        In accordo con la pianificazione temporale del progetto e lo stato attuale di avanzamento, la definizione puntuale dei casi di test per questa categoria verrà eseguita in seguito alla fase di RTB.

        \subsection{Test di Integrazione}

        In accordo con la pianificazione temporale del progetto e lo stato attuale di avanzamento, la definizione puntuale dei casi di test per questa categoria verrà eseguita in seguito alla fase di RTB.

        \subsection{Test di Sistema}

        I test di Sistema hanno lo scopo di validare il prodotto software, verificandone la conformità rispetto ai requisiti funzionali e prestazionali specificati.
        Questi test fanno riferimento ai requisiti riportati nel documento di \textit{Analisi dei Requisiti} (\S4).

            \begin{longtable}{|M{1.2cm}|M{6.3cm}|M{2.5cm}|M{1.5cm}|}
                
                \caption{Test di Sistema} \\
                \hline
                \rowcolor{gray!40}
                \textbf{ID} & \textbf{Descrizione} & \textbf{Requisito} & \textbf{Stato} \\
                \hline
                \endhead

                TS01 & Verificare che il Sistema permetta all'utente di selezionare un DIP da caricare & R1-F-Ob & NI \\
                \hline

                TS02 & Verificare che il Sistema mostri un messaggio di errore se il caricamento del DIP non avviene correttamente & R2-F-Ob & NI \\
                \hline

                TS03 & Verificare che l'utente possa rimuovere un DIP dal Sistema & R3-F-Ob & NI \\
                \hline

                TS04 & Verificare che l'utente sia in grado di visualizzare il contenuto e le informazioni del DIP & R4-F-Ob & NI \\
                \hline

                TS05 & Verificare che l'utente possa visualizzare i documenti e la relativa lista degli allegati & R5-F-Ob & NI \\
                \hline

                TS06 & Verificare che l'utente possa visualizzare ciascun allegato di ogni documento & R6-F-Ob & NI \\
                \hline

                TS07 & Verificare che l'utente possa visualizzare i dettagli di un documento & R7-F-Ob & NI \\
                \hline

                TS08 & Verificare che l'utente possa visualizzare la tipologia documentale di un documento & R8-F-Ob & NI \\
                \hline

                TS09 & Verificare che l'utente possa visualizzare il numero di un documento & R9-F-Ob & NI \\
                \hline

                TS10 & Verificare che l'utente possa visualizzare il codice registo di un documento & R10-F-Ob & NI \\
                \hline

                TS11 & Verificare che l'utente possa visualizzare la lista dei soggetti coinvolti in un documento & R11-F-Ob & NI \\
                \hline

                TS12 & Verificare che l'utente possa visualizzare un singolo soggetto coinvolto in un documento & R12-F-Ob & NI \\
                \hline

                TS13 & Verificare che l'utente possa visualizzare il ruolo di un singolo soggetto coinvolto in un documento & R13-F-Ob & NI \\
                \hline

                TS14 & Verificare che l'utente possa visualizzare l'oggetto di un documento & R14-F-Ob & NI \\
                \hline

                TS15 & Verificare che l'utente possa visualizzare il tipo aggregazione di un documento & R15-F-Ob & NI \\
                \hline

                TS16 & Verificare che l'utente possa visualizzare le informazioni di conservazione di un DIP & R16-F-Ob & NI \\
                \hline

                TS17 & Verificare che il Sistema esegua i controlli sulla validità e lo stato del processo di conservazione & R17-F-Ob & NI \\
                \hline

                TS18 & Verificare che l'utente possa selezionare "Impronta crittografica" come criterio di ricerca & R18-F-De & NI \\
                \hline

                TS19 & Verificare che l'utente possa selezionare "Identificativo documento" come criterio di ricerca & R19-F-De & NI \\
                \hline

                TS20 & Verificare che l'utente possa selezionare "Modalità di formazione" come criterio di ricerca & R20-F-De & NI \\
                \hline

                TS21 & Verificare che l'utente possa selezionare "Tipologia documentale" come criterio di ricerca & R21-F-Ob & NI \\
                \hline

                TS22 & Verificare che l'utente possa selezionare "Tipologia di flusso" come criterio di ricerca & R22-F-De & NI \\
                \hline

                TS23 & Verificare che l'utente possa selezionare "Tipo registro" come criterio di ricerca & R23-F-De & NI \\
                \hline
               
                TS24 & Verificare che l'utente possa selezionare "Data registrazione" come criterio di ricerca & R24-F-Ob & NI \\
                \hline
                
                TS25 & Verificare che l'utente possa selezionare "Ora registrazione" come criterio di ricerca & R25-F-De & NI \\
                \hline
                
                TS26 & Verificare che l'utente possa selezionare "Numero documento" come criterio di ricerca & R26-F-Ob & NI \\
                \hline

                TS27 & Verificare che l'utente possa selezionare "Codice registro" come criterio di ricerca & R27-F-Ob & NI \\
                \hline

                TS28 & Verificare che l'utente possa selezionare "Ruolo" come criterio di ricerca & R28-F-De & NI \\
                \hline

                TS29 & Verificare che l'utente possa selezionare "Tipo soggetto" come criterio di ricerca & R29-F-De & NI \\
                \hline

                TS30 & Verificare che l'utente possa selezionare "Cognome" come criterio di ricerca & R30-F-De & NI \\
                \hline

                TS31 & Verificare che l'utente possa selezionare "Nome" come criterio di ricerca & R31-F-De & NI \\
                \hline

                TS32 & Verificare che l'utente possa selezionare "Codice Fiscale" come criterio di ricerca & R32-F-De & NI \\
                \hline

                TS33 & Verificare che l'utente possa selezionare "Indirizzo digitale" come criterio di ricerca & R33-F-De & NI \\
                \hline

                TS34 & Verificare che l'utente possa selezionare "Denominazione Organizzazione" come criterio di ricerca & R34-F-De & NI \\
                \hline

                TS35 & Verificare che l'utente possa selezionare "Partita IVA" come criterio di ricerca & R35-F-De & NI \\
                \hline

                TS36 & Verificare che l'utente possa selezionare "Denominazione ufficio" come criterio di ricerca & R36-F-De & NI \\
                \hline

                TS37 & Verificare che l'utente possa selezionare "Denominazione amministrazione" come criterio di ricerca & R37-F-De & NI \\
                \hline

                TS38 & Verificare che l'utente possa selezionare "Codice IPA" come criterio di ricerca & R38-F-De & NI \\
                \hline

                TS39 & Verificare che l'utente possa selezionare "Denominazione amministrazione AOO" come criterio di ricerca & R39-F-De & NI \\
                \hline

                TS40 & Verificare che l'utente possa selezionare "Codice IPA AOO" come criterio di ricerca & R40-F-De & NI \\
                \hline

                TS41 & Verificare che l'utente possa selezionare "Denominazione amministrazione UOR" come criterio di ricerca & R41-F-De & NI \\
                \hline

                TS42 & Verificare che l'utente possa selezionare "Codice IPA UOR" come criterio di ricerca & R42-F-De & NI \\
                \hline

                TS43 & Verificare che l'utente possa selezionare "Denominazione sistema" come criterio di ricerca & R43-F-De & NI \\
                \hline

                TS44 & Verificare che l'utente possa selezionare "Oggetto" come criterio di ricerca & R44-F-Ob & NI \\
                \hline

                TS45 & Verificare che l'utente possa selezionare "Parola chiave" come criterio di ricerca & R45-F-De & NI \\
                \hline

                TS46 & Verificare che l'utente possa selezionare "Numero allegati" come criterio di ricerca & R46-F-De & NI \\
                \hline

                TS47 & Verificare che l'utente possa selezionare "Impronta crittografica allegato" come criterio di ricerca & R47-F-De & NI \\
                \hline

                TS48 & Verificare che l'utente possa selezionare "Identificativo allegato" come criterio di ricerca & R48-F-De & NI \\
                \hline

                TS49 & Verificare che l'utente possa selezionare "Descrizione allegato" come criterio di ricerca & R49-F-De & NI \\
                \hline

                TS50 & Verificare che l'utente possa selezionare "Indice di classificazione" come criterio di ricerca & R50-F-De & NI \\
                \hline

                TS51 & Verificare che l'utente possa selezionare "Descrizione dell'indice di classificazione" come criterio di ricerca & R51-F-De & NI \\
                \hline

                TS52 & Verificare che l'utente possa selezionare "Piano di classificazione" come criterio di ricerca & R52-F-De & NI \\
                \hline

                TS53 & Verificare che l'utente possa selezionare "Riservato" come criterio di ricerca & R53-F-De & NI \\
                \hline

                TS54 & Verificare che l'utente possa selezionare "Formato" come criterio di ricerca & R54-F-Ob & NI \\
                \hline

                TS55 & Verificare che l'utente possa selezionare "Nome prodotto software" come criterio di ricerca & R55-F-De & NI \\
                \hline

                TS56 & Verificare che l'utente possa selezionare "Versione prodotto software" come criterio di ricerca & R56-F-De & NI \\
                \hline

                TS57 & Verificare che l'utente possa selezionare "Produttore software" come criterio di ricerca & R57-F-De & NI \\
                \hline

                TS58 & Verificare che l'utente possa selezionare "Firmato digitalmente" come criterio di ricerca & R58-F-De & NI \\
                \hline

                TS59 & Verificare che l'utente possa selezionare "Sigillato elettronicamente" come criterio di ricerca & R59-F-De & NI \\
                \hline

                TS60 & Verificare che l'utente possa selezionare "Marcatura temporale" come criterio di ricerca & R60-F-De & NI \\
                \hline

                TS61 & Verificare che l'utente possa selezionare "Conformità coppie immagine su supporto informatico" come criterio di ricerca & R61-F-De & NI \\
                \hline

                TS62 & Verificare che l'utente possa selezionare "Tipo aggregazione" come criterio di ricerca & R62-F-Ob & NI \\
                \hline

                TS63 & Verificare che l'utente possa selezionare "Identificativo aggregazione" come criterio di ricerca & R63-F-De & NI \\
                \hline

                TS64 & Verificare che l'utente possa selezionare "Tipologia fascicolo" come criterio di ricerca & R64-F-De & NI \\
                \hline

                TS65 & Verificare che l'utente possa selezionare "Nome del documento" come criterio di ricerca & R65-F-Ob & NI \\
                \hline

                TS66 & Verificare che l'utente possa selezionare "Versione del documento" come criterio di ricerca & R66-F-De & NI \\
                \hline

                TS67 & Verificare che l'utente possa selezionare "Tipo modifica" come criterio di ricerca & R67-F-De & NI \\
                \hline

                TS68 & Verificare che l'utente possa selezionare "Data modifica" come criterio di ricerca & R68-F-De & NI \\
                \hline

                TS69 & Verificare che l'utente possa selezionare "Ora modifica" come criterio di ricerca & R69-F-De & NI \\
                \hline

                TS70 & Verificare che l'utente possa selezionare "Tempo di conservazione" come criterio di ricerca & R70-F-De & NI \\
                \hline

                TS71 & Verificare che l'utente possa selezionare "Note" come criterio di ricerca & R71-F-De & NI \\
                \hline

                TS72 & Verificare che il Sistema permetta di effettuare una ricerca impostando uno o più criteri di ricerca & R72-F-Ob & NI \\
                \hline

                TS73 & Verificare che il Sistema mostri come risultato della ricerca i documenti che soddisfano i criteri di ricerca impostati & R73-F-Ob & NI \\
                \hline

                TS74 & Verificare che l'utente possa visualizzare le informazioni relative ai documenti risultanti dalla ricerca & R74-F-Ob & NI \\
                \hline

                TS75 & Verificare che il Sistema mostri un messaggio informativo se il risultato della ricerca è vuoto & R75-F-Ob & NI \\
                \hline

                TS76 & Verificare che l'utente possa visualizzare l'anteprima di un documento selezionato & R76-F-Ob & NI \\
                \hline

                TS77 & Verificare che il Sistema mostri un messaggio di avviso se non è possibile visualizzare l'anteprima di un documento selezionato & R77-F-Ob & NI \\
                \hline

                TS78 & Verificare che l'utente possa esportare un documento selezionato, specificando il percorso dove salvarlo & R78-F-Ob & NI \\
                \hline

                TS79 & Verificare che il Sistema mostri un messaggio di errore se l'esportazione fallisce & R79-F-Ob & NI \\
                \hline

                TS80 & Verificare che l'utente possa configurare la ricerca semantica inserendo una query in linguaggio naturale & R80-F-Op & NI \\
                \hline

                TS81 & Verificare che il Sistema mostri i risultati della ricerca semantica configurata & R81-F-Op & NI \\
                \hline

                TS82 & Verificare che il Sistema permetta la stampa di un documento selezionato & R82-F-Op & NI \\
                \hline

                TS83 & Verificare che il Sistema mostri un messaggio di errore se la stampa fallisce & R83-F-Ob & NI \\
                \hline

                TS84 & Verificare che l'utente possa verificare la firma digitale di un documento selezionato & R84-F-Op & NI \\
                \hline

                TS85 & Verificare che l'utente veda il risultato a seguito della verifica della firma digitale & R85-F-Op & NI \\
                \hline

            \end{longtable}

        \subsection{Test di Accettazione}
        
            \begin{longtable}{|M{1.5cm}|M{8.5cm}|M{2cm}|}

                \caption{Test di Accettazione} \\
                \hline
                \rowcolor{gray!40}
                \textbf{ID} & \textbf{Descrizione} & \textbf{Stato} \\
                \hline
                \endhead

                TA01 & Verificare che l'utente possa caricare un pacchetto e che in seguito possa visionare la dashboard principale & NI \\
                \hline
                
                TA02 & Verificare che l'utente possa chiudere un pacchetto e che il Sistema torni alla schermata di caricamento & NI \\
                \hline

                TA03 & Verificare che l’utente possa scorrere la lista dei documenti contenuti nel pacchetto e che possa accedere ai dettagli di un singolo documento & NI \\
                \hline
                
                TA04 & Verificare che sia possibile selezionare il criterio di ricerca “Impronta crittografica” e che venga effettuata la ricerca correttamente & NI \\
                \hline

                TA05 & Verificare che, dopo la configurazione di uno o più criteri, venga mostrata all'utente la lista di documenti del pacchetto filtrata, e nel caso di ricerca senza corrispondenze venga mostrato un messaggio di avviso & NI \\
                \hline
                
                TA06 & Verificare che il prodotto dia la possibilità di effettuare una ricerca semantica tramite una query in linguaggio naturale & NI \\
                \hline

                TA07 & Verificare che venga mostrata all’utente l’anteprima di un documento dopo aver selezionato un documento e la funzionalità di anteprima & NI \\
                \hline
                
                TA08 & Verificare che l’utente possa esportare con successo un documento a seguito della selezione e dell’inserimento del percorso di esportazione & NI \\
                \hline

                TA09 & Verificare che il prodotto dia la possibilità all’utente di stampare un documento & NI \\
                \hline

                TA10 & Verificare che l’utente possa selezionare un documento firmato digitalmente e avviare la procedura di verifica, la quale, al termine dell’elaborazione, restituisce l’esito della validazione & NI \\
                \hline

            \end{longtable}


        \subsection{Checklist}

        In questa sezione sono riportate le Liste di Controllo (Checklist) utilizzate dai Verificatori durante le attività di ispezione. Queste liste costituiscono lo strumento operativo principale per accertare la conformità dei documenti agli standard di qualità prefissati dal gruppo. 
        Garantiscono uniformità stilistica, correttezza formale e coerenza interna della documentazione prodotta.

            \subsubsection{Documentazione}

                \begin{longtable}{|M{3.5cm}|M{8.5cm}|}

                    \caption{Checklist documentazione} \\
                    \hline
                    \rowcolor{gray!40}
                    \textbf{Nome} & \textbf{Descrizione} \\
                    \hline
                    \endhead

                    Immagini e tabelle & Tutte le immagini e le tabelle devono essere numerate e corredate da un titolo \\
                    \hline

                    Riferimento a risorse online & Dato che le risorse online mutano continuamente è opportuno indicare accanto al link della risorsa la data di ultima consultazione \\
                    \hline

                    Termine del glossario & Ogni volta che in un documento prodotto dal gruppo compare un termine descritto nel glossario deve essere segnalato con l’apposito stile deciso, parola\textsubscript{G}.\\
                    \hline

                    Formato delle date & Tutte le date che compaiono nel nome di un documento o nella tabella di versionamento devono avere il formato AAAA-MM-GG \\
                    \hline

                    Sezioni vuote & Tutte le sezioni dei documenti prodotti devono essere compilate \\
                    \hline

                    Cambiamento di versione & Ad ogni modifica al documento prodotta deve corrispondere un aggiornamento di versione e una riga nella tabella di versionamento \\
                    \hline

                    Errori di sintassi & Non devono essere presenti errori di sintassi \\
                    \hline

                    Forma impersonale & La forma impersonale dei verbi nei verbali è da preferire \\
                    \hline

                    Acronimi in maiuscolo & Gli acronimi devono essere in maiuscolo (es. DIP) \\
                    \hline

                    Formattazione titoli & I titoli dei documenti devono avere la prima lettera maiuscola \\
                    \hline

                    Denominazione titoli dei verbali & Il titolo dei Verbali deve essere così composto: Verbale Tipologia (Interno/Esterno) Data (nel formato DD/MM/AAAA) \\
                    \hline

                    Ordine dei nomi dei componenti & I nomi dei componenti del gruppo, all’interno dei documenti prodotti, devono essere riportati in ordine alfabetico \\
                    \hline

                    Ruoli in maiuscolo & Tutti i nomi dei ruoli devono avere la prima lettera maiuscola \\
                    \hline

                    Decisioni dei verbali & Tutte le decisioni riportate nell’apposita sezione dei verbali devono essere corredate da un identificativo univoco \\
                    \hline

                \end{longtable}

            \subsubsection{Analisi dei Requisiti}

                \begin{longtable}{|M{3.5cm}|M{8.5cm}|}
                    
                    \caption{Checklist Analisi dei requisiti} \\
                    \hline
                    \rowcolor{gray!40}
                    \textbf{Nome} & \textbf{Descrizione} \\
                    \hline
                    \endhead

                    Relazione tra casi d’uso e requisiti & Per ogni requisito deve corrispondere almeno un caso d’uso \\
                    \hline

                    Numerazione dei casi d’uso & I casi d’uso devono essere numerati in ordine crescente e nel caso di sotto-casi viene preso il numero del caso d’uso principale e viene aggiunta accanto un’altra cifra partendo da 1 \\
                    \hline

                    Diagrammi UML & Estensioni, inclusioni e generalizzazioni di un caso d’uso devono essere rappresentati in unico diagramma \\
                    \hline

                \end{longtable}
        

\newpage
\section{Cruscotto di Qualità}

In questa sezione vengono riportati gli andamenti durante tutti gli sprint delle principali metriche di qualità di processo.

    \subsection{MPC01 - Metriche Soddisfatte}
    
    \begin{figure}[H]
        \centering
        \includegraphics[width=0.9\textwidth]{quality_metrics/MPC01.png}
        \caption{Andamento percentuale delle metriche soddisfatte}
        \label{fig:mpc01}
    \end{figure}

    \noindent
    La fase di avvio riporta valori inferiori alla soglia di accettazione, questo è dovuto al carico di studio e l'instabilità dei requisiti, fattori che hanno penalizzato delle metriche specifiche.
    Superato il periodo e chiariti i dubbi progettuali, l'indice ha registrato un netto recupero a partire dallo sprint 3. Nelle fasi successive il trend si è stabilizzato su valori ottimali, raggiungendo la soglia ideale in vista della consegna della RTB.
    Anche dopo la consegna i valori rimangono ottimali, ad eccezione dello sprint 10 (\S5.2.10.7) in cui sono state apportare modifiche all'Analisi dei Requisiti su suggerimento del prof. Cardin.

    \subsection{MPC02/MPC03 - Earned Value (EV) e Planned Value (PV)}
    
    \begin{figure}[H]
        \centering
        \includegraphics[width=0.9\textwidth]{quality_metrics/MPC02.png}
        \caption{Andamento Earned Value (EV) rispetto a Planned Value (PV)}
        \label{fig:mpc02}
    \end{figure}
    
    \noindent
    Il grafico mostra un andamento allineato tra EV e PV per i primi 3 sprint. A partire dallo sprint 4 si nota una leggera flessione dell’EV dovuta ai rallentamenti del periodo festivo (come descritto nel \textit{Piano di Progetto} (\S5.2.4.7)). 
    Nello sprint 8 si registra un ulteriore incremento dello scostamento negativo tra EV e PV, recuperato poi all'interno dello sprint successivo con la conclusione del PoC. 
    Tuttavia lo scostamento si è mantenuto entro i limiti di tolleranza accettabili definiti dalla metrica MPC08. 
    In seguitop alla consegna RTB l'andamentoè rimasto generalmennte allineato grazie alle migliorie organizzative stabilite durante lo sprint 11 (\S5.3.1.7)).

    \subsection{MPC04/MPC05 - Actual Cost (AC) e Estimated To Complete (ETC)}
    
    \begin{figure}[H]
        \centering
        \includegraphics[width=0.9\textwidth]{quality_metrics/MPC04_05.png}
        \caption{Andamento di Actual Cost (AC) rispetto all'Estimated To Complete (ETC)}
        \label{fig:mpc04_05}
    \end{figure}
    
    \noindent
    La metrica analizza il costo effettivo sostenuto alla data corrente e l’ETC, cioè il budget stimato per il completamento. 
    Per la natura di queste metriche, l’andamento ottimale previsto è l'aumento progressivo di AC e la diminuzione dell’ETC, fino alla loro convergenza.
    Il gruppo è riuscito a mantenere questo andamento in modo costante.

    \subsection{MPC06 - Estimated at Completion (EAC)}
    
    \begin{figure}[H]
        \centering
        \includegraphics[width=0.9\textwidth]{quality_metrics/MPC06.png}
        \caption{Andamento della stima del costo finale del progetto (EAC)}
        \label{fig:mpc06}
    \end{figure}
    
    \noindent
    L’andamento della metrica EAC (Estimated at Completion), che rappresenta la somma tra costo sostenuto alla data di controllo e costo previsto fino alla consegna, si è mantenuto in linea con il BAC (Budget at Completion, valore totale previsto). 

    \subsection{MPC10 - Requirements Stability Index (RSI)}
    
    \begin{figure}[H]
        \centering
        \includegraphics[width=0.9\textwidth]{quality_metrics/MPC10.png}
        \caption{Andamento dell'Indice di stabilità dei requisiti (RSI)}
        \label{fig:mpc10}
    \end{figure}
    
    \noindent
    Nelle fasi iniziali il grafico mostra valori sotto la soglia di accettabilità, dovuto all’instabilità del primo periodo. Il calo registrato nello sprint 3 è dovuto al colloquio con il prof. Cardin che ha chiarito i dubbi sulla granularità dei casi d’uso, portando ad una necessaria revisione dell’ \textit{Analisi dei Requisiti}. 
    Dopo questo intervento correttivo, e grazie al costante confronto con il proponente, l’indice mostra un miglioramento, così da superare la soglia di accettabilità. Nelle fasi conclusive la metrica si è assestata nella fascia ottimale, in vista dell’RTB, in particolare dallo sprint 6 la metrica ha raggiunto il valore ideale (100\%), mantenuto anche negli sprint successivi.
    
    \subsection{MPC11 - Indice Gulpease}
    
    \begin{figure}[H]
        \centering
        \includegraphics[width=0.9\textwidth]{quality_metrics/MPC11.png}
        \caption{Andamento dell'indice di leggibilità Gulpease (valore medio dei documenti)}
        \label{fig:mpc11}
    \end{figure}
    
    \noindent
    L'andamento dell'indice mostra un decremento fisiologico nel tempo. Nelle fasi iniziali, le sezioni introduttive dei documenti hanno favorito un indice alto. Dallo sprint 3 in poi, l'aumento della complessità, l'introduzione di termini tecnici e la strutturazione formale dei requisiti hanno naturalmente abbassato la leggibilità, che si mantiene comunque stabilmente sopra la soglia di accettazione (40).

    \subsection{MPC15 - Rischi non calcolati}
    
    \begin{figure}[H]
        \centering
        \includegraphics[width=0.9\textwidth]{quality_metrics/MPC15.png}
        \caption{Andamento dell'indice dei Rischi non calcolati}
        \label{fig:mpc15}
    \end{figure}
    
    \noindent
    Il valore massimo registrato è stato 1 rischio per sprint, quindi sempre ampiamente entro la soglia di accettabilità e mai vicino al limite critico.
    La sequenza evidenzia una fase iniziale stabile (0–0), seguita da una breve area di instabilità in cui si concentrano tre sprint con valore 1 consecutivo. Questa variazione è attribuibile a ritardi nella consegna che hanno temporaneamente aumentato l’incertezza nella gestione dei rischi non previsti.
    Successivamente, il trend ritorna a stabilità completa (0) per un numero significativo di sprint consecutivi, indicando un miglioramento nella gestione del rischio e nella prevedibilità delle attività.
    
    Nel complesso, il grafico mostra un processo sotto controllo, con una varianza molto bassa, episodi isolati e rapidamente rientrati, e nessun superamento della soglia accettabile.

    \subsection{MPC16 - Efficienza Temporale}

    \begin{figure}[H]
        \centering
        \includegraphics[width=0.9\textwidth]{quality_metrics/MPC16.png}
        \caption{Andamento dell'efficienza temporale}
        \label{fig:mpc16}
    \end{figure}
    
    \noindent
    Il grafico evidenzia l'impatto della curva di apprendimento sul lavoro produttivo. I valori iniziali sopra la soglia, con il picco massimo nello sprint 2, sono giustificati dall'intensa attività di studio individuale necessaria per comprendere il materiale fornito dall'azienda.
    Successivamente l'indice migliora nettamente, ad eccezione di una risalita nello sprint 5 dovuta alle necessità di comprendere le tecnologie per l'avvio della stesura del codice del PoC.

    \subsection{MPD04 - Code Coverage}
    
    \begin{figure}[H]
        \centering
        \includegraphics[width=0.9\textwidth]{quality_metrics/MPD04.png}
        \caption{Andamento dell'indice di Code Coverage}
        \label{fig:mpd04}
    \end{figure}
    
    \noindent
    L'indice di Code Coverage è monitorabile con l'inizio della codifica dell'MVP. Pur avendo avuto un picco iniziale, seguito da un decremento, l'indice è rimasto sempre sotto la soglia minima di accettazione. Ciò è dovuto alla scarsa quantità di test implementati.

\end{document}
